\documentclass{article}

% if you need to pass options to natbib, use, e.g.:
%     \PassOptionsToPackage{numbers, compress}{natbib}
% before loading neurips_2026

% The authors should use one of these tracks.
% Before accepting by the NeurIPS conference, select one of the options below.
% 0. "default" for submission
% \usepackage{neurips_2026}
% the "default" option is equal to the "main" option, which is used for the Main Track with double-blind reviewing.
% 1. "main" option is used for the Main Track
 \usepackage[main]{neurips_2026}
% 2. "position" option is used for the Position Paper Track
%  \usepackage[position]{neurips_2026}
% 3. "eandd" option is used for the Evaluations & Datasets Track
 % \usepackage[eandd]{neurips_2026}
 % if you want to opt-in for a de-anomymized submission (not recommended, but OK):
 %\usepackage[eandd, nonanonymous]{neurips_2026}
% 4. "creativeai" option is used for the Creative AI Track
%  \usepackage[creativeai]{neurips_2026}
% 5. "sglblindworkshop" option is used for the Workshop with single-blind reviewing
 % \usepackage[sglblindworkshop]{neurips_2026}
% 6. "dblblindworkshop" option is used for the Workshop with double-blind reviewing
%  \usepackage[dblblindworkshop]{neurips_2026}

% After being accepted, the authors should add "final" behind the track to compile a camera-ready version.
% 1. Main Track
 % \usepackage[main, final]{neurips_2026}
% 2. Position Paper Track
%  \usepackage[position, final]{neurips_2026}
% 3. Evaluations & Datasets Track
 % \usepackage[eandd, final]{neurips_2026}
% 4. Creative AI Track
%  \usepackage[creativeai, final]{neurips_2026}
% 5. Workshop with single-blind reviewing
%  \usepackage[sglblindworkshop, final]{neurips_2026}
% 6. Workshop with double-blind reviewing
%  \usepackage[dblblindworkshop, final]{neurips_2026}
% Note. For the workshop paper template, both \title{} and \workshoptitle{} are required, with the former indicating the paper title shown in the title and the latter indicating the workshop title displayed in the footnote.
% For workshops (5., 6.), the authors should add the name of the workshop, "\workshoptitle" command is used to set the workshop title.
% \workshoptitle{WORKSHOP TITLE}

% "preprint" option is used for arXiv or other preprint submissions
 % \usepackage[preprint]{neurips_2026}

% to avoid loading the natbib package, add option nonatbib:
%    \usepackage[nonatbib]{neurips_2026}

\usepackage[utf8]{inputenc} % allow utf-8 input
\usepackage[T1]{fontenc}    % use 8-bit T1 fonts
\usepackage{hyperref}       % hyperlinks
\usepackage{url}            % simple URL typesetting
\usepackage{booktabs}       % professional-quality tables
\usepackage{amsfonts}       % blackboard math symbols
\usepackage{amsmath}        % math environments
\usepackage{amssymb}        % math symbols
\usepackage{nicefrac}       % compact symbols for 1/2, etc.
\usepackage{microtype}      % microtypography
\usepackage{xcolor}         % colors
\usepackage{graphicx}       % figures
\usepackage{multirow}       % table multirow
\usepackage{bm}             % bold math

% Note. For the workshop paper template, both \title{} and \workshoptitle{} are required, with the former indicating the paper title shown in the title and the latter indicating the workshop title displayed in the footnote. 
\title{Formatting Instructions For NeurIPS 2026}


% The \author macro works with any number of authors. There are two commands
% used to separate the names and addresses of multiple authors: \And and \AND.
%
% Using \And between authors leaves it to LaTeX to determine where to break the
% lines. Using \AND forces a line break at that point. So, if LaTeX puts 3 of 4
% authors names on the first line, and the last on the second line, try using
% \AND instead of \And before the third author name.


\author{%
  Muhammad Umar Farooq\\
  College of Innovation and Technology\\
  University of Michigan\\
  Flint, MI 48502 \\
  \texttt{mufarooq@umich.edu} \\
  % examples of more authors
  \And
  Khalid Malik \\
  College of Innovation and Technology\\
  University of Michigan\\
  Flint, MI 48502 \\
  \texttt{drmalik@umich.edu} \\
  % \AND
  % Coauthor \\
  % Affiliation \\
  % Address \\
  % \texttt{email} \\
  % \And
  % Coauthor \\
  % Affiliation \\
  % Address \\
  % \texttt{email} \\
  % \And
  % Coauthor \\
  % Affiliation \\
  % Address \\
  % \texttt{email} \\
}


\begin{document}


\maketitle


\begin{abstract}
TODO: Write abstract.
\end{abstract}

%% TODO: Section 1 — Introduction
%% TODO: Section 2 — Related Work

% ======================================================================
\section{Proposed method}
\label{sec:method}
% ======================================================================

We present \textbf{NeSyDeFake}, a neuro-symbolic framework for generalizable deepfake detection that fuses neural perception with symbolic reasoning under a principled uncertainty-aware objective.
The key insight is that face manipulations inevitably violate semantic consistency constraints---relationships between facial attributes that hold for authentic faces but break under synthesis.
Rather than relying solely on learned visual features that overfit to source-specific artifacts, NeSyDeFake explicitly encodes and verifies these constraints, producing interpretable violation signals that generalize across manipulation methods.

The framework produces three independent evidence streams---\emph{spatial}, \emph{concept}, and \emph{causal}---that are fused via Evidential Deep Learning (EDL) into a Dirichlet distribution over classes, yielding both a prediction and a calibrated epistemic uncertainty estimate.
We describe each component below and summarize the full pipeline in Figure~\ref{fig:architecture}.

\begin{figure}[t]
  \centering
  \fbox{\rule[-.5cm]{0cm}{4cm} \rule[-.5cm]{12cm}{0cm}}
  \caption{Overview of the NeSyDeFake framework. A frozen CLIP backbone extracts spatial features that produce spatial evidence via a linear head. Precomputed semantic and forensic features are processed by a concept branch (hand-coded consistency rules) and a causal constraint verification branch (learned constraints, forensic anomaly detection, counterfactual prediction). All three evidence streams are fused via Confidence-Modulated Evidence Fusion (CMEF) under a novel Neuro-Symbolic Evidential Deep Learning (NeSy-EDL) objective with per-branch auxiliary supervision and inter-branch disagreement calibration.}
  \label{fig:architecture}
\end{figure}


% ------------------------------------------------------------------
\subsection{Multi-source feature extraction}
\label{sec:features}
% ------------------------------------------------------------------

NeSyDeFake operates on three precomputed feature sets extracted from face-cropped frames (scale $1.3{\times}$, capturing hair, ears, and neck context).

\paragraph{Spatial features.}
A frozen CLIP ViT-L/14 backbone extracts 1024-dimensional visual representations.
Only the LayerNorm parameters are fine-tuned during training, following the linear probing with LayerNorm tuning protocol~\citep{gend}, which preserves the pretrained feature space while allowing task-specific adaptation.

\paragraph{Semantic features.}
We construct a 122-dimensional combined semantic vector from two sources:
(i)~\emph{fast-extracted features} ($\bm{s}_\text{fast} \in \mathbb{R}^{58}$) computed by InsightFace, MediaPipe, and DeepFace at ${\sim}30$ms/frame, covering demographics (continuous gender, age, ethnicity scores), expression probabilities (8 emotions), action unit activations (17 FACS AUs), head pose (pitch, yaw, roll), gaze direction (6D), and landmark geometry (eye/mouth aspect ratios, jaw symmetry); and
(ii)~\emph{vision-language model features} ($\bm{s}_\text{vlm} \in \mathbb{R}^{64}$) from a FaceBench-finetuned Face-LLaVA 13B model, capturing attributes that require holistic visual reasoning (accessories, hair style, skin texture, lighting, background).
The combined vector is $\bm{s} = [\bm{s}_\text{fast} \| \bm{s}_\text{vlm}] \in \mathbb{R}^{122}$.

\paragraph{Forensic features.}
An 83-dimensional forensic vector $\bm{f} \in \mathbb{R}^{83}$ captures pixel-level manipulation traces across five semantic groups:
\emph{boundary/texture} (12 features: Sobel gradients and Laplacian blur at face parsing boundaries),
\emph{symmetry/color} (18: left-right intensity asymmetry and cross-region color histogram distances, DCT frequency energy, quality metrics),
\emph{patch noise} (16: pairwise and cross-channel noise consistency between facial regions),
\emph{SRM noise} (27: Steganalysis Rich Model residuals---5 kernels $\times$ mean, standard deviation, kurtosis, plus multi-scale noise statistics),
and \emph{FFT spectral} (10: radial spectral energy in 8 frequency bins, spectral slope, and high-frequency ratio).


% ------------------------------------------------------------------
\subsection{Spatial evidence branch}
\label{sec:spatial}
% ------------------------------------------------------------------

The spatial branch provides the base detection signal.
The frozen CLIP features $\bm{h} \in \mathbb{R}^{1024}$ pass through a learnable projection head $g_\phi$ followed by a linear classification head:
\begin{equation}
  \bm{e}_\text{spatial} = \operatorname{softplus}\bigl(W_\text{cls}\, g_\phi(\bm{h}) + \bm{b}_\text{cls}\bigr) \in \mathbb{R}^{2}_{\geq 0},
  \label{eq:spatial_evidence}
\end{equation}
where the softplus activation ensures non-negative evidence suitable for the Dirichlet framework (Section~\ref{sec:edl}).
This branch alone corresponds to a standard CLIP-based detector and serves as the foundation upon which the neuro-symbolic components build.


% ------------------------------------------------------------------
\subsection{Concept evidence branch}
\label{sec:concept}
% ------------------------------------------------------------------

The concept branch encodes domain knowledge about facial attribute relationships as differentiable \emph{consistency rules}.
We define 23 rules organized into five categories, each producing a scalar violation score $v_j \in [0, 1]$:

\paragraph{Mutual exclusion rules} (3 rules) detect contradictory attributes:
\begin{equation}
  v_\text{mouth} = \sigma(\text{mouth\_open}) \cdot \sigma(\text{mouth\_closed}),
  \label{eq:mutual_excl}
\end{equation}
and analogously for lighting-conflict and symmetry-conflict pairs. For authentic faces, at most one attribute in each pair should be active, so $v_j \approx 0$; manipulated faces that paste features from different sources often violate these mutual exclusions.

\paragraph{Expression--action unit coherence rules} (9 rules) enforce the Facial Action Coding System.
Each emotion should co-occur with specific AUs; violations indicate that the expression was transferred without the corresponding muscle activations:
\begin{equation}
  v_\text{happy-AU6} = \bigl|\sigma(\text{expr\_happy}) - \sigma(\text{AU6})\bigr|.
  \label{eq:au_rule}
\end{equation}

\paragraph{Cross-region identity rules} (5 rules) capture face/context inconsistencies visible in the $1.3{\times}$ crop, e.g., $v_\text{gender-beard} = \sigma(-\text{gender\_score}) \cdot \sigma(\text{beard})$, which fires when a face classified as female has a visible beard---a common artifact in identity-swapping deepfakes.

\paragraph{Skin-age and structural rules} (6 rules) check physiological consistency (wrinkles should correlate with age, left-right gaze directions should agree, facial landmarks should be bilaterally symmetric).

The 23-dimensional violation vector $\bm{v} = [v_1, \ldots, v_{23}]^\top$ is concatenated with the combined semantic features and mapped to evidence:
\begin{equation}
  \bm{e}_\text{concept} = \operatorname{softplus}\bigl(\text{MLP}_\text{concept}([\bm{s} \| \bm{v}])\bigr) \in \mathbb{R}^{2}_{\geq 0},
  \label{eq:concept_evidence}
\end{equation}
where $\text{MLP}_\text{concept}$ is a two-layer network with LayerNorm, GELU activation, and dropout.


% ------------------------------------------------------------------
\subsection{Causal constraint verification branch}
\label{sec:ccv}
% ------------------------------------------------------------------

While the 23 hand-coded rules capture known attribute interactions, they cannot cover all consistency patterns that distinguish real from manipulated faces.
The \emph{Causal Constraint Verification} (CCV) branch extends symbolic reasoning with three learned components that discover additional violation signals from data.

\subsubsection{Learned constraint functions}
\label{sec:lcf}

We learn $K{=}16$ additional soft constraint functions over the combined semantic space.
A shared trunk $t(\bm{s}) = \operatorname{GELU}(W_t\,\operatorname{LN}(\bm{s}) + \bm{b}_t)$ maps the 122-dimensional input to a hidden representation, from which $K$ independent sigmoid heads produce violation scores:
\begin{equation}
  \bm{c} = \sigma(W_h\, t(\bm{s}) + \bm{b}_h) \in [0, 1]^K.
  \label{eq:learned_constraints}
\end{equation}
Each learned constraint discovers an attribute interaction pattern that the hand-coded rules miss.
The shared trunk amortizes computation while the independent heads specialize, analogous to multi-task learning with a shared encoder.
Weights are initialized near zero so that all constraints start at $\bm{c} \approx 0.5$ (maximum uncertainty), allowing the model to learn which interactions are discriminative.


\subsubsection{Forensic anomaly detector}
\label{sec:fad}

Pixel-level forensic features of authentic faces exhibit consistent patterns within semantic regions (e.g., uniform noise characteristics across skin patches), while manipulated faces show anomalous patterns at blending boundaries.
Rather than learning a single monolithic model over all 83 features, we train five lightweight autoencoders, one per forensic group $g \in \{\text{boundary}, \text{symmetry}, \text{noise}_\text{patch}, \text{noise}_\text{srm}, \text{spectral}\}$:
\begin{equation}
  a_g = \frac{1}{d_g}\bigl\|\bm{f}_g - \operatorname{dec}_g(\operatorname{enc}_g(\bm{f}_g))\bigr\|_2^2,
  \label{eq:anomaly}
\end{equation}
where $\bm{f}_g \in \mathbb{R}^{d_g}$ is the forensic feature subvector for group $g$, the encoder maps to an 8-dimensional bottleneck with LayerNorm and GELU, and the decoder reconstructs.
Real faces, which exhibit internally consistent forensic patterns, are well-reconstructed ($a_g \approx 0$); manipulated faces that introduce inconsistent noise or boundary artifacts produce high reconstruction errors.
The five per-group anomaly scores $\bm{a} = [a_1, \ldots, a_5]^\top$ provide interpretable, semantically grouped forensic signals.


\subsubsection{Counterfactual consistency predictor}
\label{sec:counterfactual}

The counterfactual predictor poses the question: \emph{``If a face truly possessed these semantic and forensic attributes, what would it look like in the visual feature space?''}
For authentic faces, the predicted appearance matches the actual CLIP representation; for manipulated faces, the attributes are inconsistent with the visual content, producing a mismatch.

Concretely, we compress the CLIP spatial features $\bm{h}$ to a compact representation $\bm{z} = W_c\, \bm{h} \in \mathbb{R}^{32}$ (with the gradient detached from the backbone).
A predictor network maps the full semantic context $[\bm{s} \| \bm{v} \| \bm{f}] \in \mathbb{R}^{228}$ to a predicted appearance vector:
\begin{equation}
  \hat{\bm{z}} = \text{MLP}_\text{cf}\bigl([\bm{s} \| \bm{v} \| \bm{f}]\bigr) \in \mathbb{R}^{32},
  \quad\quad
  m = \|\bm{z} - \hat{\bm{z}}\|_2^2,
  \label{eq:counterfactual}
\end{equation}
where $m \in \mathbb{R}_{\geq 0}$ is the counterfactual mismatch score.
This formulation directly operationalizes the causal hypothesis: authentic faces are generated by a consistent process where attributes determine appearance, while manipulations break this mapping.


\subsubsection{Evidence fusion}

The CCV branch concatenates all verification signals and maps them to evidence:
\begin{equation}
  \bm{e}_\text{causal} = \operatorname{softplus}\bigl(\text{MLP}_\text{ccv}([\bm{v} \| \bm{c} \| \bm{a} \| m])\bigr) \in \mathbb{R}^2_{\geq 0},
  \label{eq:ccv_evidence}
\end{equation}
where the input is a 45-dimensional vector (23 hand-coded violations + 16 learned constraints + 5 anomaly scores + 1 counterfactual mismatch).
The MLP has LayerNorm, GELU, and dropout.
This entire branch contains approximately 15K trainable parameters, making it lightweight relative to the backbone.


% ------------------------------------------------------------------
\subsection{Neuro-symbolic evidential deep learning (NeSy-EDL)}
\label{sec:edl}
% ------------------------------------------------------------------

A central challenge in multi-branch neuro-symbolic detection is \emph{how to fuse heterogeneous evidence sources under uncertainty}.
Standard EDL~\citep{sensoy2018evidential} assumes a single evidence vector; naively summing evidence from neural and symbolic branches ignores that their reliability varies per sample and fails to reflect reasoning \emph{conflict} in the fused uncertainty.
We introduce \textbf{NeSy-EDL}, extending EDL with three novel components that make the evidential framework multi-branch-aware.


\subsubsection{Confidence-modulated evidence fusion (CMEF)}
\label{sec:cmef}

Standard gated fusion uses static scalar gates: $\bm{e}_\text{total} = \bm{e}_\text{spatial} + \sigma(\gamma_1)\bm{e}_\text{concept} + \sigma(\gamma_2)\bm{e}_\text{causal}$, where the contribution of each branch is identical across all samples.
This is suboptimal: a symbolic branch may be highly informative for one sample (e.g., clear rule violations) but uninformative for another (e.g., ambiguous attributes).

CMEF replaces static gates with \emph{per-sample dynamic weighting} based on each branch's own Dirichlet strength:
\begin{equation}
  \bm{e}_\text{total} = \bm{e}_\text{spatial}
    + \sigma(\gamma_1)\,\phi(S_\text{concept})\,\bm{e}_\text{concept}
    + \sigma(\gamma_2)\,\phi(S_\text{causal})\,\bm{e}_\text{causal},
  \label{eq:cmef}
\end{equation}
where $S_b = \sum_k (\bm{e}_b + \bm{1})_k$ is the Dirichlet strength of branch $b$, and the \emph{confidence gate} is:
\begin{equation}
  \phi(S_b) = \sigma\!\left(\frac{S_b - K}{\tau}\right), \quad \tau = \exp(\theta_\tau),
  \label{eq:conf_gate}
\end{equation}
with learnable log-temperature $\theta_\tau$ (initialized to 0, so $\tau{=}1$).
When $S_b = K$ (uniform Dirichlet, maximum epistemic uncertainty), $\phi \approx 0.5$; when $S_b \gg K$ (confident branch), $\phi \to 1$; when $S_b < K$ (degenerate), $\phi < 0.5$.
The learnable gates $\gamma_1, \gamma_2$ are initialized to $-0.5$ and $-0.8$, respectively.

\paragraph{Key property.}
On out-of-distribution data where the neural spatial branch becomes uncertain (low $S_\text{spatial}$), the symbolic branches---whose rule-based and constraint-based evidence generalizes across manipulation methods---often retain higher confidence.
CMEF automatically allocates greater relative weight to these symbolic branches, enabling \emph{dynamic neural-symbolic trust allocation} precisely in the generalization regime that matters.


\subsubsection{Dirichlet prediction and uncertainty}

The fused evidence parameterizes a Dirichlet distribution:
\begin{equation}
  \bm{\alpha} = \bm{e}_\text{total} + \bm{1}, \quad
  S = \textstyle\sum_{k=1}^{K}\alpha_k, \quad
  \hat{p}_k = \frac{\alpha_k}{S}, \quad
  u = \frac{K}{S},
  \label{eq:dirichlet}
\end{equation}
where $\bm{\alpha} \in \mathbb{R}^K_{>0}$ are Dirichlet concentration parameters ($K{=}2$), $\hat{p}_k$ is the expected class probability, and $u \in (0, 1]$ is the epistemic uncertainty.


\subsubsection{Training objective}
\label{sec:loss}

The NeSy-EDL objective combines five terms, the first three inherited from standard EDL and the last two novel to our multi-branch setting.

\paragraph{Evidential log-likelihood.}
Type-II maximum likelihood under the Dirichlet prior on the \emph{fused} evidence:
\begin{equation}
  \mathcal{L}_\text{nll} = -\sum_{k=1}^{K} y_k\bigl(\psi(\alpha_k) - \psi(S)\bigr),
  \label{eq:edl_nll}
\end{equation}
where $\bm{y}$ is the one-hot target and $\psi$ is the digamma function.
Per-sample class weights $[w_\text{real}{=}1.8,\, w_\text{fake}{=}0.6]$ counteract CLIP's bias toward the fake class.

\paragraph{Annealed KL divergence.}
A KL penalty between the predicted Dirichlet (with correct-class evidence removed) and the uniform prior $\text{Dir}(\bm{1})$:
\begin{equation}
  \mathcal{L}_\text{kl} = \text{KL}\bigl[\text{Dir}(\tilde{\bm{\alpha}}) \| \text{Dir}(\bm{1})\bigr], \quad
  \tilde{\alpha}_k = (1 - y_k)(\alpha_k - 1) + 1,
  \label{eq:kl}
\end{equation}
annealed from 0 to $\lambda_\text{kl}{=}0.15$ over the first 10 epochs to prevent premature evidence suppression.

\paragraph{Accuracy-vs-uncertainty calibration (AVU).}
Following~\citet{karandikar2021soft}:
\begin{equation}
  \mathcal{L}_\text{avu} = \bm{1}_\text{correct}\bigl[-\hat{p}_\text{max}\log(1{-}u)\bigr]
    + \bm{1}_\text{wrong}\bigl[-(1{-}\hat{p}_\text{max})\log(u)\bigr],
  \label{eq:avu}
\end{equation}
which penalizes confident-but-wrong predictions and rewards uncertainty on misclassified samples.

\paragraph{Per-branch auxiliary supervision (PBAS).}
\emph{Novel.}
In multi-branch fusion, a dominant branch (here, the CLIP spatial branch with 1024-d features) can suppress gradient flow to weaker symbolic branches, causing them to learn noise rather than meaningful evidence.
PBAS applies a lightweight standalone EDL loss to \emph{each branch independently}:
\begin{equation}
  \mathcal{L}_\text{aux} = \sum_{b \in \mathcal{B}} \mathcal{L}_\text{edl}^{(b)}\bigl(\bm{e}_b,\, \bm{y}\bigr),
  \label{eq:pbas}
\end{equation}
where $\mathcal{B} = \{\text{spatial}, \text{concept}, \text{causal}\}$ and $\mathcal{L}_\text{edl}^{(b)}$ is the standard EDL loss (NLL + annealed KL) applied to branch $b$'s evidence alone.
This ensures each branch produces independently meaningful evidence before fusion, analogous to auxiliary losses in multi-task learning but operating in the evidential space.

\paragraph{Inter-branch disagreement calibration (IBDC).}
\emph{Novel.}
Standard EDL has no mechanism for the fused uncertainty to reflect reasoning \emph{conflict} between evidence sources---it only reflects total evidence magnitude.
IBDC introduces a calibration loss that ties epistemic uncertainty to inter-branch agreement:
\begin{equation}
  d(\bm{x}) = 1 - \cos\bigl(\hat{\bm{p}}_\text{neural}(\bm{x}),\, \hat{\bm{p}}_\text{symbolic}(\bm{x})\bigr),
  \label{eq:disagreement}
\end{equation}
\begin{equation}
  \mathcal{L}_\text{bdc} = -d\,\log(u) - (1{-}d)\,\log(1{-}u),
  \label{eq:ibdc}
\end{equation}
where $\hat{\bm{p}}_b = \bm{\alpha}_b / S_b$ is the Dirichlet mean of branch $b$, $d \in [0, 1]$ is the cosine disagreement (detached from the computation graph), and $u$ is the fused epistemic uncertainty.
This binary cross-entropy loss encourages the model to be uncertain when branches \emph{disagree} (e.g., CLIP says real but rules detect violations) and confident when they \emph{agree}.
The disagreement is averaged over all $\binom{|\mathcal{B}|}{2}$ branch pairs.

\paragraph{Total objective.}
\begin{equation}
  \mathcal{L} = \mathcal{L}_\text{nll}
    + \lambda_\text{kl}\,\beta(t)\,\mathcal{L}_\text{kl}
    + \lambda_\text{avu}\,\mathcal{L}_\text{avu}
    + \lambda_\text{aux}\,\mathcal{L}_\text{aux}
    + \lambda_\text{bdc}\,\mathcal{L}_\text{bdc},
  \label{eq:total_loss}
\end{equation}
where $\beta(t) = \min(1, t/T_\text{anneal})$ is the annealing schedule, $\lambda_\text{kl}{=}0.15$, $\lambda_\text{avu}{=}0.1$, $\lambda_\text{aux}{=}0.1$, and $\lambda_\text{bdc}{=}0.05$.
The CMEF parameters ($\theta_\tau$, $\gamma_1$, $\gamma_2$) are trained jointly with the detection objective, receiving gradients from all five loss terms.


% ------------------------------------------------------------------
\subsection{Source-paired training}
\label{sec:paired}
% ------------------------------------------------------------------

Following~\citet{gend}, we train with source-paired batches where each real face is paired with its corresponding manipulated version from the same source video.
This ensures that the model cannot exploit dataset-level distributional differences between real and fake sets and must instead learn manipulation-specific signals.
Combined with class-weighted EDL and the neuro-symbolic consistency verification, this training protocol encourages the model to learn generalizable detection criteria rather than source-specific shortcuts.

\begin{table}[t]
  \caption{Summary of evidence sources and their dimensionality. The NeSyDeFake framework fuses three evidence streams, each capturing complementary signals for deepfake detection.}
  \label{tab:evidence_summary}
  \centering
  \small
  \begin{tabular}{llrl}
    \toprule
    \textbf{Branch} & \textbf{Input signals} & \textbf{Dim} & \textbf{Nature} \\
    \midrule
    Spatial & CLIP ViT-L/14 features & 1024 & Neural (learned visual) \\
    \midrule
    \multirow{2}{*}{Concept} & Combined semantic features & 122 & \multirow{2}{*}{Symbolic (rule-based)} \\
     & Hand-coded consistency rules & 23 & \\
    \midrule
    \multirow{4}{*}{Causal (CCV)} & Hand-coded rules (shared) & 23 & Symbolic \\
     & Learned constraint functions & 16 & Neuro-symbolic \\
     & Forensic anomaly scores & 5 & Neural (autoencoder) \\
     & Counterfactual mismatch & 1 & Neuro-symbolic \\
    \midrule
    \textbf{Fused evidence} & CMEF Dirichlet aggregation & 2 & NeSy-EDL \\
    \bottomrule
  \end{tabular}
\end{table}

\begin{ack}
Use unnumbered first level headings for the acknowledgments. All acknowledgments
go at the end of the paper before the list of references. Moreover, you are required to declare
funding (financial activities supporting the submitted work) and competing interests (related financial activities outside the submitted work).
More information about this disclosure can be found at: \url{https://neurips.cc/Conferences/2026/PaperInformation/FundingDisclosure}.


Do {\bf not} include this section in the anonymized submission, only in the final paper. You can use the \texttt{ack} environment provided in the style file to automatically hide this section in the anonymized submission.
\end{ack}

\section*{References}


References follow the acknowledgments in the camera-ready paper. Use unnumbered first-level heading for
the references. Any choice of citation style is acceptable as long as you are
consistent. It is permissible to reduce the font size to \verb+small+ (9 point)
when listing the references.
Note that the Reference section does not count towards the page limit.
\medskip


{
\small


[1] Alexander, J.A.\ \& Mozer, M.C.\ (1995) Template-based algorithms for
connectionist rule extraction. In G.\ Tesauro, D.S.\ Touretzky and T.K.\ Leen
(eds.), {\it Advances in Neural Information Processing Systems 7},
pp.\ 609--616. Cambridge, MA: MIT Press.


[2] Bower, J.M.\ \& Beeman, D.\ (1995) {\it The Book of GENESIS: Exploring
  Realistic Neural Models with the GEneral NEural SImulation System.}  New York:
TELOS/Springer--Verlag.


[3] Hasselmo, M.E., Schnell, E.\ \& Barkai, E.\ (1995) Dynamics of learning and
recall at excitatory recurrent synapses and cholinergic modulation in rat
hippocampal region CA3. {\it Journal of Neuroscience} {\bf 15}(7):5249-5262.
}


%%%%%%%%%%%%%%%%%%%%%%%%%%%%%%%%%%%%%%%%%%%%%%%%%%%%%%%%%%%%

\appendix

\section{Technical appendices and supplementary material}
Technical appendices with additional results, figures, graphs, and proofs may be submitted with the paper submission before the full submission deadline (see above). You can upload a ZIP file for videos or code, but do not upload a separate PDF file for the appendix. There is no page limit for the technical appendices. 

Note: Think of the appendix as ``optional reading'' for reviewers. The paper must be able to stand alone without the appendix; for example, adding critical experiments that support the main claims to an appendix is inappropriate. 

%%%%%%%%%%%%%%%%%%%%%%%%%%%%%%%%%%%%%%%%%%%%%%%%%%%%%%%%%%%%

\newpage
\input{checklist.tex}


\end{document}